\noindent\null\hfill{\small\textbf{ZUR HANNOVERSCHEN TRIANGULATION.}}\hfill\makebox[0pt][r]{\fontsize{13}{16}\selectfont \textbf{375}}\par\medskip

\begin{center}
\textsc{Gauss} an \textsc{Olbers}.\quad Göttingen, 25.~Februar 1825.
\end{center}

\par\noindent\hspace*{\parindent}\mbox{.\ .\ .\ .\ .\ .}\quad Jetzt noch ein paar Worte über das geodätische Nivellement.
Ich bin zwar selbst mit mir über diesen Gegenstand \textls[120]{ganz auf dem Reinen},
allein ich weiss nicht, ob ich mich in der Kürze so darüber werde erklären
können, dass ich Sie sofort zur Übereinstimmung bringen werde.

\par\noindent\hspace*{\parindent}Ich habe immer geglaubt, dass der Ausdruck »Localattraction« sehr
übel gewählt ist und leicht verkehrte Ansichten veranlassen kann. Man
sollte sagen, dass die Richtungen der Schwere nicht mit dem Gange, der bei
einem gleichförmigen Sphäroid stattfinden würde, Schritt halten. Die
Richtung der Schwere ist das Totalproduct der Anziehung aller Bestandtheile des
Erdkörpers (und der Centrifugalkraft), und bei dessen unregelmässiger
Zusammensetzung in Rücksicht der Dichtigkeit, sowie bei den Unebenheiten auf
der Oberfläche wird jene nicht dieselbe sein können, wie bei einem
regelmässigen Sphäroid. Allein \textls[120]{wie auch die Zusammensetzung sei}, immer
wird durch jeden Punkt eine Fläche, die ganz um die Erde herum geht,
gelegt werden können, auf welcher die Richtung der Schwere \textls[120]{genau} senkrecht
ist, und die Oberfläche einer zusammenhängenden ruhigen Flüssigkeit würde
dieselbe vorstellen. \textbf{Diese} Fläche ist es, die eine Horizontalfläche heisst (couche
de niveau); den Punkten \textls[120]{dieser} Fläche \textls[120]{legt man gleiche Höhe bei}, ohne
sich im mindesten darum zu bekümmern, ob oder wie viel sie von einem
elliptischen Sphäroid abweichen, und die Höhen über dieser Fläche gibt sowohl
das Barometer als die trigonometrische Messung an, so dass beide immer mit
einander übereinstimmen müssen. Dabei wird bloss vorausgesetzt\textsuperscript{*}), dass auf
jeder Dreieckslinie die Richtung der Schwere sich nach dem Gesetz der
Stetigkeit ändert (obgleich vielleicht schneller oder langsamer als bei dem
elliptischen Sphäroid), und diese Voraussetzung kann nur dann eine kleine
Unrichtigkeit hervorbringen, wenn an der einen Dreiecksstation eine \textls[120]{wahre}
Localattraction stattfindet, die bloss örtlich und auf einen kleinen Raum beschränkt

\par\smallskip\noindent\rule{.24\textwidth}{.4pt}\par\smallskip
{\footnotesize
\noindent\hspace*{.06\linewidth}\textsuperscript{*}) und natürlich auch, dass alle Zenithdistanzen reciprok gemessen werden, was bei meiner Messung
\textls[120]{ohne Ausnahme} gilt. Bei \textls[120]{einseitigen} Messungen ist Ihre Bemerkung vollkommen gegründet, da
man dabei die Amplitude sphäroidisch berechnen muss.\par
}

\clearpage
\noindent\makebox[0pt][l]{\fontsize{13}{16}\selectfont \textbf{376}}\hfill{\small\textbf{BRIEFWECHSEL MIT OLBERS.}}\hfill\null\par\medskip
\noindent (ausserhalb desselben unmerklich) ist. Allein ich halte mich überzeugt, dass,
den Brocken höchstens ausgenommen, eine \textls[120]{solche} Localattraction im ganzen
Umfange meiner und der \textsc{Schumachers}chen Dreiecke nicht statthat.\ \mbox{.\ .\ .\ .\ .\ .}
